\section{Numerical Guarantees for Double-Word Arithmetic over IEEE~754}
\label{sec:numerics}

To prove the numerical guarantees presented in \Cref{sec:overview},
we define \emph{functional models} of the arithmetic algorithms.
These mathematical models are expressed as Rocq functions using
IEEE~754 operations from VCFloat2 \citep{AppelKellison2024} and the
binary floating-point representation from Flocq \citep{Boldo2011}.
Each model follows the sequence of rounded operations in the
corresponding arithmetic kernel and returns a pair of floating-point
values, independently of the C representation. We define the models
over a generic radix-two format; \FenceDD{} uses the binary64
instances.

We first establish the rounding properties and error-free
transformation results needed to analyze complete double-word
operations. We then continue the division example from
\Cref{sec:overview-division}, giving a proof sketch of the numerical
guarantee and explaining how the proofs for the remaining operations
differ.

\subsection{FLX, FLT, and IEEE~754 Arithmetic}
\label{sec:flx-to-ieee}

To extend the earlier mechanization by \citeauthor{Rideau2022} to
IEEE~754 arithmetic \citep{IEEE754-2019}, we relate rounding in two
Flocq formats with the same radix and precision: floating-point with
unbounded exponents (FLX) and floating-point with gradual underflow
(FLT). The earlier mechanization uses FLX and therefore excludes
both underflow and overflow. FLT imposes a lower exponent bound and
includes subnormal values.

For radix $\beta>1$, precision $p>0$, and lower exponent $e_{\min}$, the two
formats represent
\[
\begin{aligned}
  \mathbb{F}_{\mathrm{FLX}}(\beta,p)
    &=
    \left\{m\beta^q \mathrel{}\middle|\mathrel{}
      m,q\in\mathbb{Z},\ |m|<\beta^p\right\},\\
  \mathbb{F}_{\mathrm{FLT}}(\beta,p,e_{\min})
    &=
    \left\{m\beta^q \mathrel{}\middle|\mathrel{}
      m,q\in\mathbb{Z},\ |m|<\beta^p,\ q\geq e_{\min}\right\}.
\end{aligned}
\]
Thus,
$\mathbb{F}_{\mathrm{FLT}}(\beta,p,e_{\min})
 \subseteq\mathbb{F}_{\mathrm{FLX}}(\beta,p)$.

The lower exponent bound determines the smallest positive
representable value $\eta$ and the smallest positive normal value $\nu$:
\begin{equation}
  \eta=\beta^{e_{\min}}
  \qquad\text{and}\qquad
  \nu=\beta^{e_{\min}+p-1}.
  \label{eq:flt-thresholds}
\end{equation}
Nonzero representable values with magnitude below $\nu$ are
subnormal. In this range, adjacent representable values are spaced
$\eta$ apart. For binary64, $\eta=2^{-1074}$ and $\nu=2^{-1022}$.

Let $\operatorname{RN}_{\mathrm{FLX}}$ and
$\operatorname{RN}_{\mathrm{FLT}}$ denote rounding to nearest with
ties to even in the two formats. For exact real values with
magnitude at least $\nu$, the Flocq theorem
\texttt{FLT.round\_FLT\_FLX} gives
\begin{equation}
  |x|\geq\nu
  \quad\Longrightarrow\quad
  \operatorname{RN}_{\mathrm{FLT}}(x)
    =\operatorname{RN}_{\mathrm{FLX}}(x).
  \label{eq:round-flt-flx}
\end{equation}
This equality allows us to reuse the earlier FLX arguments for
operations whose exact results have magnitude at least $\nu$.
For smaller exact results, the two roundings may differ, so we
analyze FLT rounding directly.

Neither FLX nor FLT has an upper exponent bound. Our error analysis
assumes that the returned components are finite and uses this
hypothesis to relate the IEEE~754 computations to FLT rounding.
We separately establish sufficient conditions under which the
returned components are finite.

\subsection{Error-Free Transformations over IEEE~754}
\label{sec:error-free-transformations}

Double-word algorithms use error-free transformations that return
a rounded sum or product together with a correction intended to
equal the exact rounding error. To extend the analysis of these
algorithms to IEEE~754 arithmetic, we must determine whether this
correction remains exact when gradual underflow occurs. For
addition, the equality remains valid under suitable conditions.
For multiplication, the exact rounding error may be too small to
represent. We identify a sufficient condition under which the
correction remains exact and otherwise bound the error in the sum
of the returned components.

\begin{definition}[Error-free transformation]
\label[definition]{def:error-free-transformation}
For an arithmetic operation with exact real result $r$, an error-free
transformation returns a pair $(r_h,r_\ell)$ satisfying
\[
  r_h=\operatorname{RN}(r)
  \qquad\text{and}\qquad
  r=r_h+r_\ell.
\]
\end{definition}

Thus, $r_\ell=r-r_h$ is the exact rounding error. Every pair returned by an
error-free transformation also forms a double word because
$r_h=\operatorname{RN}(r_h+r_\ell)$ follows from
\Cref{def:error-free-transformation}.

\Cref{fig:error-free-transformations} gives the three error-free
transformations used by the double-word algorithms in \FenceDD{}. 
We use the naming conventions and algorithms from 
\citet[Algorithms\ 1--3]{Joldes2017}.
\TwoSum{} is due to Knuth, while
\FastTwoSum{} and the original error-free transformations for products are due
to Dekker \citep{Knuth1997,Dekker1971}.

\begin{figure}[t]
  \centering
  \begingroup
  \setlength{\fboxsep}{6pt}
  \setlength{\fboxrule}{0.4pt}
  \begin{subfigure}[t]{\columnwidth}
    \phantomsubcaption\label{alg:two-sum}
    \fcolorbox{gray!35}{gray!3}{%
      \begin{minipage}[t]{\dimexpr\linewidth-2\fboxsep-2\fboxrule\relax}
        \centering
        \textbf{(\thesubfigure) \TwoSum{}$(a,b)$}
        \par\smallskip
        \raggedright\small
        \textit{Input:} Floating-point values $a$ and $b$\par
        \textit{Output:} A rounded sum $s$ and correction $t$
        \setlength{\abovedisplayskip}{4pt}
        \setlength{\belowdisplayskip}{0pt}
        \[
          \begin{array}{@{}r@{\hspace{0.6em}}l@{}}
            1: & s \leftarrow \operatorname{RN}(a+b)\\
            2: & a' \leftarrow \operatorname{RN}(s-b)\\
            3: & b' \leftarrow \operatorname{RN}(s-a')\\
            4: & \delta_a \leftarrow \operatorname{RN}(a-a')\\
            5: & \delta_b \leftarrow \operatorname{RN}(b-b')\\
            6: & t \leftarrow \operatorname{RN}(\delta_a+\delta_b)\\
            7: & \mathbf{return}\ (s,t)
          \end{array}
        \]
      \end{minipage}}
  \end{subfigure}

  \vspace{7pt}

  \begin{subfigure}[t]{\columnwidth}
    \phantomsubcaption\label{alg:fast-two-sum}
    \fcolorbox{gray!35}{gray!3}{%
      \begin{minipage}[t]{\dimexpr\linewidth-2\fboxsep-2\fboxrule\relax}
        \centering
        \textbf{(\thesubfigure) \FastTwoSum{}$(a,b)$}
        \par\smallskip
        \raggedright\small
        \textit{Input:} Floating-point values $a$ and $b$\par
        \textit{Output:} A rounded sum $s$ and correction $t$
        \setlength{\abovedisplayskip}{4pt}
        \setlength{\belowdisplayskip}{0pt}
        \[
          \begin{array}{@{}r@{\hspace{0.6em}}l@{}}
            1: & s \leftarrow \operatorname{RN}(a+b)\\
            2: & z \leftarrow \operatorname{RN}(s-a)\\
            3: & t \leftarrow \operatorname{RN}(b-z)\\
            4: & \mathbf{return}\ (s,t)
          \end{array}
        \]
      \end{minipage}}
  \end{subfigure}

  \vspace{7pt}

  \begin{subfigure}[t]{\columnwidth}
    \phantomsubcaption\label{alg:fast-two-mult}
    \fcolorbox{gray!35}{gray!3}{%
      \begin{minipage}[t]{\dimexpr\linewidth-2\fboxsep-2\fboxrule\relax}
        \centering
        \textbf{(\thesubfigure) \FastTwoMult{}$(a,b)$}
        \par\smallskip
        \raggedright\small
        \textit{Input:} Floating-point values $a$ and $b$\par
        \textit{Output:} A rounded product $\pi_h$ and correction $\pi_\ell$
        \setlength{\abovedisplayskip}{4pt}
        \setlength{\belowdisplayskip}{0pt}
        \[
          \begin{array}{@{}r@{\hspace{0.6em}}l@{}}
            1: & \pi_h \leftarrow \operatorname{RN}(ab)\\
            2: & \pi_\ell \leftarrow
                 \operatorname{RN}(ab-\pi_h)
                 \quad\text{(fused multiply-add)}\\
            3: & \mathbf{return}\ (\pi_h,\pi_\ell)
          \end{array}
        \]
      \end{minipage}}
  \end{subfigure}
  \endgroup
\caption{Error-free transformations used by the double-word algorithms
implemented in \FenceDD.
\TwoSum{} and \FastTwoSum{} return a rounded sum and correction;
\FastTwoMult{} returns a rounded product and FMA correction.
\Cref{thm:addition-transformations,thm:fast-two-sum-small,thm:fast-two-mult}
state when each correction is exact.}
  \label{fig:error-free-transformations}
\end{figure}

\subsubsection{Exact corrections for addition.}

Gradual underflow does not affect the corrections
returned by \TwoSum{} and \FastTwoSum{}. When the returned components are
finite, \TwoSum{} returns the rounded sum and its exact rounding error.
\FastTwoSum{} also returns the rounded sum and its exact rounding error
under the usual ordering condition on the operands. 
These results, proved by \citet{Boldo2017}, allow the later
double-word proofs to use the addition transformations without imposing a
lower bound on the magnitude of the operands.

\begin{theorem}[Error-free transformations for addition over IEEE~754]
\label{thm:addition-transformations}
Let $a$ and $b$ be finite values in a radix-two format, and suppose the
returned components are finite.
\begin{enumerate}
  \item If $(s,t)=\TwoSum{}(a,b)$, then
    $s=\operatorname{RN}(a+b)$ and $s+t=a+b$.
  \item If $|b|\leq|a|$ and $(s,t)=\FastTwoSum{}(a,b)$, then $(s,t)$
    satisfies $s=\operatorname{RN}(a+b)$ and $s+t=a+b$.
\end{enumerate}
\end{theorem}

The ordering condition $|b|\leq|a|$ is sufficient but not necessary
for exactness, and it cannot be established for every use of
\FastTwoSum{} in the double-word algorithms. We prove another
sufficient condition needed by the later analysis: assuming finite
returned components, the correction remains exact when
$|a+b|\leq2\nu$.

\begin{theorem}[\FastTwoSum{} for small exact sums]
\label{thm:fast-two-sum-small}
Let $\nu$ be the smallest positive normal value in the format, as defined in
\Cref{eq:flt-thresholds}.  If $(s,t)=\FastTwoSum{}(a,b)$ has finite components
and $|a+b|\leq2\nu$, then $s=\operatorname{RN}(a+b)$ and $s+t=a+b$.
\end{theorem}

The proof appears in Appendix~\ref{app:fast-two-sum-small-proof}; like all 
results in this section, it is checked in Rocq.

\subsubsection{Corrections for multiplication near underflow.}

\FastTwoMult{} returns
$\pi_h=\operatorname{RN}(ab)$ and
$\pi_\ell=\operatorname{RN}(d)$, where $d=ab-\pi_h$ is the exact rounding
error. If $d$ is representable, then $\pi_\ell=d$ and the returned components
sum to $ab$; near underflow, however, $d$ may be too small to represent.
Theorem~2 of \citet[Section~3]{BoldoDaumas2003} proves representability when
$a=n_a2^{e_a}$ and $b=n_b2^{e_b}$ satisfy
$e_a+e_b\geq e_{\min}$. From this result, we derive a sufficient condition on $|ab|$ under which the
returned components sum to the exact product. When this condition does not
hold but both components are finite, we instead bound the difference between
their sum and the exact product.

\begin{theorem}[\FastTwoMult{} over IEEE~754]
\label{thm:fast-two-mult}
Let $a$ and $b$ be finite values in a radix-two format with precision
$p\geq3$ and lower exponent $e_{\min}$. Define
\[
  \eta\triangleq2^{e_{\min}}
  \qquad\text{and}\qquad
  \tau\triangleq2^{e_{\min}+2p-1}.
\]
Suppose $(\pi_h,\pi_\ell)=\FastTwoMult{}(a,b)$ returns two finite components.
Then
\[
  \pi_h=\operatorname{RN}(ab)
  \qquad\text{and}\qquad
  \left|(\pi_h+\pi_\ell)-ab\right|\leq\frac{\eta}{2}.
\]
If $ab=0$ or $|ab|\geq\tau$, then
\[
  \pi_\ell=ab-\pi_h
  \qquad\text{and}\qquad
  \pi_h+\pi_\ell=ab.
\]
\end{theorem}

\begin{proof}[Proof sketch]
Let $d=ab-\pi_h$ be the exact rounding error. Lines~1 and~2 of
\Cref{alg:fast-two-mult} give $\pi_h=\operatorname{RN}(ab)$ and
$\pi_\ell=\operatorname{RN}(d)$, so
$|(\pi_h+\pi_\ell)-ab|=|\pi_\ell-d|$.

If $ab=0$, then $d=\pi_\ell=0$. If $|ab|\geq\tau$, representations
$a=n_a2^{e_a}$ and $b=n_b2^{e_b}$ cannot satisfy
$e_a+e_b<e_{\min}$, since that inequality would imply $|ab|<\tau$.
Theorem~2 of \citet[Section~3]{BoldoDaumas2003} therefore shows that $d$ is
representable, giving $\pi_\ell=d$ and exact reconstruction.

In the remaining case, $0<|ab|<\tau$. The FLT rounding bound gives
$|d|<2\nu$, where $\nu=2^{e_{\min}+p-1}$. Representable values in this range
are spaced by $\eta$, so $|\pi_\ell-d|\leq\eta/2$, which proves the stated
error bound.
\end{proof}

\subsection{Numerical Guarantee for Division of a Double Word by a
Floating-Point Value}
\label{sec:division-numerics}

Division of a double word by a floating-point value shows how the preceding
error-free transformation results combine within a complete double-word
operation. We identify sufficient conditions under which the returned pair
forms a double word and its represented value satisfies an error bound with
relative and absolute terms. 
The proof accounts for both effects of the finite exponent
range: the relevant computations must remain finite, and operations involving
subnormal values may require direct analysis in FLT.

Given a double word $x=x_h+x_\ell$ and a nonzero floating-point divisor $y$,
the exact quotient is
\begin{equation}
  q=\frac{x_h+x_\ell}{y}.
  \label{eq:dw-div-exact-quotient}
\end{equation}
The functional model \DWDivFP{} represents the rounded operations used to
approximate this quotient. \Cref{alg:dw-div-fp} gives its mathematical
definition. Each $\operatorname{RN}$ denotes rounding to nearest with ties to even.

\begin{figure}[t]
  \centering
  \begingroup
  \setlength{\fboxsep}{6pt}
  \setlength{\fboxrule}{0.4pt}
  \fcolorbox{gray!35}{gray!3}{%
    \begin{minipage}[t]{\dimexpr\linewidth-2\fboxsep-2\fboxrule\relax}
      \centering
      \textbf{\DWDivFP{}$(x_h,x_\ell,y)$}
      \par\smallskip
      \raggedright\small
      \textit{Input:} A double word $(x_h,x_\ell)$ and a nonzero
      floating-point value $y$\par
      \textit{Output:} A pair $(z_h,z_\ell)$ approximating $q$
      \setlength{\abovedisplayskip}{4pt}
      \setlength{\belowdisplayskip}{0pt}
      \[
        \begin{array}{@{}r@{\hspace{0.6em}}l@{}}
          1: & t_h \leftarrow \operatorname{RN}(x_h/y)\\
          2: & (\pi_h,\pi_\ell) \leftarrow \FastTwoMult{}(t_h,y)\\
          3: & \delta_h \leftarrow \operatorname{RN}(x_h-\pi_h)\\
          4: & \delta_t \leftarrow \operatorname{RN}(\delta_h-\pi_\ell)\\
          5: & \delta \leftarrow \operatorname{RN}(\delta_t+x_\ell)\\
          6: & t_\ell \leftarrow \operatorname{RN}(\delta/y)\\
          7: & (z_h,z_\ell) \leftarrow \FastTwoSum{}(t_h,t_\ell)\\
          8: & \mathbf{return}\ (z_h,z_\ell)
        \end{array}
      \]
    \end{minipage}}
  \endgroup
  \caption{The functional model \DWDivFP{} for dividing a double word by a
  floating-point value.}
  \label{alg:dw-div-fp}
\end{figure}

The model follows the identity
\[
  q=\frac{x}{y}
    =t_h+\frac{x-t_hy}{y}.
\]
Line~1 computes the leading quotient $t_h$.
Lines~2--5 approximate the residual $x-t_hy$, including the contribution of
the low word $x_\ell$. Line~6 divides this residual by $y$ to obtain the
correction $t_\ell$. Line~7 constructs the final pair.

\subsubsection{Sufficient Conditions for the Numerical Guarantee}
\label{sec:division-conditions}

We state sufficient conditions for \DWDivFP{} to return a finite
double word satisfying an error bound. We state the operand
conditions separately from the assumption that the returned
components are finite, then give stronger conditions under which
result finiteness follows.

For operands $x_h,x_\ell,y$, let
$(z_h,z_\ell)=\DWDivFP{}(x_h,x_\ell,y)$, and let $t_h$ be the rounded leading
quotient computed on Line~1 of \Cref{alg:dw-div-fp}. When $y\neq0$, write
$q_0=x_h/y$ for the exact leading quotient. We write
$\operatorname{Fin}(v_1,\ldots,v_n)$ when all its arguments are finite. Let
$\Omega$ be the least power of two greater than every finite magnitude in the
format; for binary64, $\Omega=2^{1024}$.

\begin{definition}[Conditions for \DWDivFP{}]
\label[definition]{def:dw-div-conditions}
Define
\begin{equation}
\begin{aligned}
\mathcal{I}_{\mathrm{div}}(x_h,x_\ell,y)
  \quad\triangleq\quad&
  (x_h,x_\ell)\text{ is a double word}\\
  &{}\land
  \left(x_h=0
    \;\lor\;
    2^{e_{\min}+2p}\leq |x_h|\right),\\[2pt]
\mathcal{F}_{\mathrm{div}}(x_h,x_\ell,y)
  \quad\triangleq\quad&
  \operatorname{Fin}(z_h,z_\ell),\\[2pt]
\mathcal{S}_{\mathrm{div}}(x_h,x_\ell,y)
  \quad\triangleq\quad&
  \mathcal{I}_{\mathrm{div}}(x_h,x_\ell,y)
  \land \operatorname{Fin}(x_h,x_\ell,y,t_h)\\
  &{}\land |x_h|\leq\Omega/4
  \land |q_0|\leq\Omega/4.
\end{aligned}
\label{eq:dw-div-conditions}
\end{equation}
\end{definition}

The condition $\mathcal{I}_{\mathrm{div}}$ requires the dividend to be a
double word and imposes a lower bound on its high word unless the dividend is
zero. It places no lower bound on $q_0$, so the leading quotient may be
subnormal. The condition $\mathcal{F}_{\mathrm{div}}$ states that both
returned words are finite, as required for the error bound and double-word
relation. The stronger condition $\mathcal{S}_{\mathrm{div}}$ adds finiteness
and upper bounds sufficient to derive
$\mathcal{F}_{\mathrm{div}}$ without assuming it separately.

\begin{lemma}[Finiteness from the Stronger Condition]
\label[lemma]{lem:dw-div-safe-implies-finite}
If $\mathcal{S}_{\mathrm{div}}(x_h,x_\ell,y)$ holds, then
$\mathcal{F}_{\mathrm{div}}(x_h,x_\ell,y)$ holds.
\end{lemma}

\subsubsection{Error Bounds and Double-Word Result}

The main theorem assumes
$\mathcal{I}_{\mathrm{div}}$ and $\mathcal{F}_{\mathrm{div}}$. Under these
conditions, the functional model returns a double word and satisfies an error
bound with relative and absolute terms. Under additional magnitude conditions,
the absolute term is unnecessary. By the preceding finiteness lemma, the same
conclusions also follow from $\mathcal{S}_{\mathrm{div}}$.

\begin{theorem}[Numerical guarantee for \DWDivFP{}]
\label{thm:dw-div-fp-numerical}
Let the underlying format be radix two with precision $p\geq4$, lower
exponent $e_{\min}$, unit roundoff $u=2^{-p}$, and
$\eta=2^{e_{\min}}$. Let $\nu=2^{e_{\min}+p-1}$ be its smallest positive
normal value, and let
$(z_h,z_\ell)=\DWDivFP{}(x_h,x_\ell,y)$.
Suppose
$\mathcal{I}_{\mathrm{div}}(x_h,x_\ell,y)$ and
$\mathcal{F}_{\mathrm{div}}(x_h,x_\ell,y)$ hold. These conditions imply
$y\neq0$. Let $q$ be the exact quotient defined in
\Cref{eq:dw-div-exact-quotient}, and let $q_0=x_h/y$ be the exact leading
quotient. Then $(z_h,z_\ell)$ is a double word:
\[
  z_h=\operatorname{RN}(z_h+z_\ell).
\]
Its represented value satisfies
\begin{equation}
  \left|(z_h+z_\ell)-q\right|
    \leq 3u^2|q|+\eta.
  \label{eq:dw-div-fp-mixed}
\end{equation}
Let $\delta$ denote the residual computed on Line~5 of
\Cref{alg:dw-div-fp}. If $|q_0|\geq\nu$ and either $\delta=0$ or
$|\delta/y|\geq\nu$, the sharper bound holds:
\begin{equation}
  \left|(z_h+z_\ell)-q\right|
    \leq 3u^2|q|.
  \label{eq:dw-div-fp-relative}
\end{equation}
\end{theorem}

\begin{proof}[Proof sketch]
The conditions imply that $y\neq0$ and that the relevant intermediate results
are finite. If $x_h=0$, the double-word condition gives $x_\ell=0$, and the
result follows directly. Otherwise,
$|x_h|\geq2^{e_{\min}+2p}$.

First suppose $|q_0|\geq\nu$. Rounding on Line~1 agrees in FLT
and FLX. The lower bound on
$|x_h|$ gives $|t_hy|\geq2^{e_{\min}+2p-1}$, so
\Cref{thm:fast-two-mult} gives $\pi_h+\pi_\ell=t_hy$. Sterbenz's lemma and
representability then carry the FLX residual analysis through Lines~3--5
\citep{Sterbenz1974}. If $\delta=0$ or $|\delta/y|\geq\nu$, Line~6 also agrees
in the two formats and gives \Cref{eq:dw-div-fp-relative}; otherwise, its FLT
and FLX results differ by at most $\eta$, giving the mixed bound. Finally,
$|t_\ell|\leq|t_h|$, so \FastTwoSum{} preserves their sum and returns a double
word by \Cref{thm:addition-transformations}.

Now suppose $|q_0|<\nu$. The fixed spacing in FLT gives
$|t_h-q_0|\leq\eta/2$. The double-word
condition gives $|x_\ell/y|<\eta/2$, and hence $|t_h-q|\leq\eta$. The lower
bound on $|x_h|$ also implies $|y|>2^{p+1}$. Applying the FLT rounding bounds
and \Cref{thm:fast-two-mult} to Lines~2--5 gives
\begin{equation}
  |\delta-(x_h+x_\ell-t_hy)|
    \leq\frac{\eta|y|}{2},
  \qquad
  \left|\frac{\delta}{y}-(q-t_h)\right|\leq\frac{\eta}{2}.
  \label{eq:dw-div-subnormal-residual}
\end{equation}
Line~6 adds at most another $\eta/2$, and $|t_h+t_\ell|\leq2\nu$.
\Cref{thm:fast-two-sum-small} therefore shows that Line~7 preserves the sum
and returns a double word, with
\[
  |(z_h+z_\ell)-q|
    \leq\eta
    \leq3u^2|q|+\eta.
\]
This proves both error bounds and the double-word property.
\end{proof}

By \Cref{lem:dw-div-safe-implies-finite},
$\mathcal{S}_{\mathrm{div}}$ implies $\mathcal{F}_{\mathrm{div}}$ and already
includes $\mathcal{I}_{\mathrm{div}}$. Consequently,
\Cref{thm:dw-div-fp-numerical} gives the following two alternatives.

\begin{corollary}[Conditions for the Division Guarantee]
\label[corollary]{cor:dw-div-conditions}
The conclusions of \Cref{thm:dw-div-fp-numerical} hold if
\[
  \mathcal{I}_{\mathrm{div}}(x_h,x_\ell,y)
  \land
  \mathcal{F}_{\mathrm{div}}(x_h,x_\ell,y)
\]
holds, or if $\mathcal{S}_{\mathrm{div}}(x_h,x_\ell,y)$ holds.
\end{corollary}

\Cref{thm:division-checks} proves that the weak and result finiteness checks in
the default wrapper establish the first alternative, while the strong check in
the alternative wrapper establishes the second.

\subsection{Numerical Guarantees for the Remaining Operations}
\label{sec:remaining-numerical-conditions}

The five remaining functional models also return double words and satisfy the
error bounds summarized in \Cref{tab:operation-preview} under
operation-specific conditions. Rather than repeat their algorithms and full
proofs, we collect the conditions needed to extend the existing analyses to
the finite IEEE~754 exponent range and state one common theorem.

Write $x=(x_h,x_\ell)$ and, where present, $y=(y_h,y_\ell)$ for the
double-word operands, and let $a$ denote a floating-point operand. For
addition and multiplication, the returned components must be finite, but no
additional condition on the operands is required beyond the double-word
assumptions. Division of two double words additionally requires
%
\begin{equation}
  2^{e_{\min}+2p}
  \leq |y_h|
  \leq \frac{\Omega}{2^{p+5}}.
  \label{eq:dw-div-dw-divisor-condition}
\end{equation}
%
In every case, finiteness may either be assumed directly for the returned
components or established from the sufficient conditions in
\Cref{tab:remaining-numerical-conditions}.
%
\begin{table}[t]
\centering
\caption{Sufficient conditions for the five remaining functional
models to return finite components. Each row provides an
alternative to the result-finiteness assumption in
\Cref{thm:remaining-numerical-guarantees}; all other hypotheses
still apply. Here $x=(x_h,x_\ell)$ and $y=(y_h,y_\ell)$ are
double-word operands, $a$ is a floating-point operand, and
$\operatorname{Fin}$ states that all listed values are finite.
The constant $\Omega$ is the least power of two greater than
every finite magnitude in the underlying format.}
\label{tab:remaining-numerical-conditions}
\setlength{\tabcolsep}{3pt}
\begin{tabular}{@{}p{0.31\columnwidth}p{0.64\columnwidth}@{}}
\toprule
Operation & Sufficient conditions for a finite result \\
\midrule
Addition $(x,a)$
  & $\operatorname{Fin}(x_h,x_\ell,a)$;
    $|x_h|,|a|\leq\Omega/4$ \\
\addlinespace
Addition $(x,y)$
  & $\operatorname{Fin}(x_h,x_\ell,y_h,y_\ell)$;
    $|x_h|,|y_h|\leq\Omega/4$ \\
\addlinespace
Multiplication $(x,a)$
  & $\operatorname{Fin}(x_\ell,\operatorname{RN}(x_h a))$;
    $|x_h a|\leq\Omega/2$ \\
\addlinespace
Multiplication $(x,y)$
  & $\operatorname{Fin}(x_\ell,y_\ell,\operatorname{RN}(x_h y_h))$;
    $|x_h y_h|\leq\Omega/2$ \\
\addlinespace
Division $(x,y)$
  & $\operatorname{Fin}(x_h,x_\ell,y_h,y_\ell)$;
    $|x_h/y_h|\leq\Omega/4$ \\
\bottomrule
\end{tabular}
\end{table}
%
Under these conditions, each functional model returns a double word and
satisfies its operation-specific error bound:
%
\begin{theorem}[Numerical guarantees for the remaining operations]
\label{thm:remaining-numerical-guarantees}
Let the underlying format be radix two with precision $p$, unit
roundoff $u=2^{-p}$, and smallest positive representable value
$\eta=2^{e_{\min}}$. Choose one of the five operations in
\Cref{tab:remaining-numerical-conditions}, and let $(z_h,z_\ell)$
be the pair returned by the corresponding functional model.
Let $r$ be the exact real result of the operation on the values
represented by the operands.

Assume $p\geq3$ for the addition operations and multiplication by a
floating-point value, $p\geq5$ for multiplication of two double
words, and $p\geq14$ for division of two double words. Suppose every
double-word operand satisfies \Cref{def:double-word}. For division
of two double words, also suppose
\Cref{eq:dw-div-dw-divisor-condition} holds.

Finally, suppose either that the returned components are finite,
$\operatorname{Fin}(z_h,z_\ell)$, or that the sufficient conditions
for finiteness in the corresponding row of
\Cref{tab:remaining-numerical-conditions} hold.

Then both returned components are finite and form a double word:
\[
  z_h=\operatorname{RN}(z_h+z_\ell).
\]
The represented value $\widehat r=z_h+z_\ell$ satisfies the
corresponding error bound on $|\widehat r-r|$ in
\Cref{tab:operation-preview}, using $u$ and $\eta$ for the underlying
format.
\end{theorem}

\paragraph{Proof overview.}
\emph{Addition.}
The error-free transformations for addition continue to return the exact
rounding error when intermediate results are subnormal by
\Cref{thm:addition-transformations}. Consequently, the existing relative
error bounds require no lower bound on the operand magnitudes.

\emph{Multiplication.}
When a product correction is too small to be represented exactly,
\Cref{thm:fast-two-mult} bounds the lost correction. Propagating that error
through the remaining computation produces the absolute terms in the 
\Cref{tab:operation-preview} bounds.

\emph{Division of two double words.}
In contrast to division by a floating-point value in
\Cref{alg:dw-div-fp}, division of two double words begins by approximating the
reciprocal of $y_h$. The lower bound in
\Cref{eq:dw-div-dw-divisor-condition} prevents this approximation from
overflowing. The upper bound implies $\eta|y_h|\leq u^2/4$, allowing the error
from a lost product correction to be absorbed into the relative term. 
The bound on the quotient in \Cref{tab:remaining-numerical-conditions} limits the
magnitude of the leading quotient and is sufficient to prove that the returned
components are finite.
